\section{Introduction}

Recent advances in large language models~\citep{gpt,gpt2,gpt3,gpt3.5,gpt4,llama,llama2} open the possibility of building language agents that can autonomously solve complex tasks. The common practice for developing AI agents is to decompose complex tasks into LLM pipelines where prompts and tools are stacked together~\citep{park2023generative,hong2023metagpt,zhou2023agents,chen2023agentverse,OpenAgents}. In a sense, language agents can be viewed as AI systems that connect connectionism AI (i.e., the LLM backbone of agents) and symbolism AI (i.e., the pipeline of prompts and tools), which partially explains their effectiveness in real-world problem-solving scenarios.

However, the current state of language agents development is limited by the extensive engineering effort required to build and customize language agent systems for a specific task. Specifically, researchers and developers have to manually decompose complex tasks into subtasks, which we refer to as nodes, that are more tractable for LLMs and then carefully design prompts and tools, including API functions, knowledge bases, memories, etc., for specific nodes. The complexity of this process makes the current landscape of language agent research \textit{model-centric}, or \textit{engineering-centric}. This means it is almost impossible for researchers to manually tune or optimize language agents on datasets on which we can train neural nets in a \textit{data-centric} way. This limits the robustness and versatility of manually coded language agents and requires substantial engineering effort to adapt language agents to new tasks or data distributions. We believe the transition from engineering-centric language agents development to data-centric learning is an important step in language agent research.

To this end, a number of recent efforts has been made on automatic optimization of language agents. For example, DSpy~\citep{khattab2023dspy} introduces a framework for algorithmically optimizing LLM prompts via bootstrapping or random searching in a combinatory space of different prompt components and GPTSwarm~\citep{zhuge2024language} further proposes to tackle the combinatorial optimization challenge raised in DSPy via an iterative optimization process. Agent-pro~\citep{zhang2024agentpro} proposes a framework to optimize the prompts components corresponding to the agents' internal policy in competitive environments. AgentOptimizer~\citep{zhang2024offline} proposes a framework to optimize functions with carefully engineered prompts. While effective in some scenarios, these approaches only optimize separate modules in an agent system such as a prompt for a specific node. As a result, these optimization methods are prone to local optimum of isolated prompts, tools, and nodes that leads to compromised performance for the entire agent system. This resembles the early practice in training neural nets~\citep{hinton2006reducing} where layers are separately optimized and it now seems trivial that optimizing neural nets as a whole leads to better performance. We believe that this is also the case in agent optimization and jointly optimization of all symbolic components within an agent is the key for optimizing agents. 

\begin{figure}[tp]
    \centering
    \includegraphics[width=0.8\linewidth]{figs/overview.pdf}
    \caption{Analogy between \baby and neural nets connectionist learning.}
    \label{fig:overview}
    \vspace{-15pt}
\end{figure}

In this work, we introduce a \textit{\baby} framework for training language agents. The \baby framework is inspired by the connectionist learning procedure~\citep{hinton1990connectionist} used for training neural nets. To be specific, we make an analogy between language agents and neural nets: the agent pipeline of an agent corresponds to the computational graph of a neural net, a node in the agent pipeline corresponds to a layer in the neural net, and the prompts and tools for a node correspond to the weights of a layer. In this way, we are able to implement the main components of connectionist learning, i.e., backward propagation and gradient-based weight update, in the context of agent training using language-based loss, gradients, and weights. Specifically, we implement loss function, back-propagation, and weight optimizer in the context of agent training with carefully designed prompt pipelines. For a training example, our framework first conducts the ``forward pass'' (agent execution) and stores the input, output, prompts, and tool usage in each node in a ``trajectory''. We then use a prompt-based loss function to evaluate the outcome, resulting in a ``language loss''. Afterward, we back-propagate the language loss from the last to the first node along the trajectory, resulting in textual analyses and reflections for the symbolic components within each node, we call them language gradients. Finally, we update all symbolic components in each node, as well as the computational graph consisting of the nodes and their connections, according to the language gradients with another carefully designed prompt. Our approach also naturally supports optimizing multi-agent systems by considering nodes as different agents or allowing multiple agents to take actions in one node.

The \baby framework is an agent learning framework that mimics the standard connectionist learning procedure. In contrast to existing methods that either optimize single prompt or tool in a separate manner, the \baby framework jointly optimizes all symbolic components within an agent system, including prompts, tools, and the pipeline that stacks them into an agent system. This top-down optimization scheme also enables the \baby framework to optimize the agent system ``holistically'', avoiding local optimum for each separated component. This makes it possible for language agents targeting complex real-world problems to effectively \textit{learn from data}, opening up the possibility to transform the current state of language agent research from engineering-centric to data-centric. Moreover, since the language-based loss function does not require ground-truth when generating the language loss, our framework enables language agents to \textit{learn from experience} and \textit{delibrately} update all their symbolic components after being created and deployed in the wild, enabling ``self-evolving agents''\footnote{Agents can also collect training data in the wild and update the LLM backbone via fine-tuning. In this way, all components in the agent can be updated. We leave this for future work.}.

As a proof-of-concept, we conduct a series of experiments on both standard LLM benchmarks and complex agentic tasks. Our results demonstrate the effectiveness of the proposed \baby framework to optimize and design prompts and tools, as well as update the overall agent pipeline by learning from training data. We open-source all codes and prompts in the \baby framework to facilitate future research on \textit{data-centric} agent learning. 

%The symbolic optimization framework can be easily wrapped into an external tool which an agent can use after being deployed in the wild.

